% =============================================================================
% Section 4 — Discussion
% =============================================================================
\section{Discussion}
\label{sec:discussion}

\subsection{Pathway partition: why bass ``feels'' different from treble}
\label{sec:why-bass-feels}

Our results partition the ``feel it in your gut'' experience into its
constituent pathways.  The airborne acoustic contribution, computed using
the energy-consistent Junger \& Feit reciprocity formulation, reaches
$\sim$\SI{0.35}{\percent} of the ISO~2631 perception threshold at peak
EDM levels ($\SI{115}{\decibel}$ at \SI{40}{\hertz}).  This is a
``near-miss''---\num{2.5}~orders of magnitude below threshold, not the
four orders suggested by the pressure-based approach---but it remains
insufficient for airborne perception alone.

The structure-borne pathway (floor and seat vibration) dominates by a
factor of ${\sim}40$ under conservative assumptions.  With realistic
venue floors (lightweight timber, resonant panels), the structure-borne
displacement can readily exceed the perception threshold, explaining the
well-documented visceral sensation at concerts.

The key insight is that the airborne pathway is not negligibly weak---it
is a significant fraction of the overall budget, particularly in
near-field conditions or at extreme SPLs.  The ``feel it in your gut''
experience is best understood as a \emph{pathway partition problem}: the
structure-borne pathway carries ${\sim}\SI{97}{\percent}$ of the
perceptible energy, with the airborne component contributing a small
but non-zero remainder.

\subsection{Additional coupling pathways}
\label{sec:additional-pathways}

Beyond the two quantified pathways, several additional mechanisms
contribute to bass perception:
%
\begin{enumerate}
  \item \textbf{Chest-wall coupling.}  The large surface area of the
    chest wall acts as a more efficient acoustic receiver than the
    abdomen.  Chest-wall motion at sub-bass frequencies can produce
    thoracic pressure fluctuations that are transmitted to the
    abdominal cavity.

  \item \textbf{Somatosensory stimulation.}  High-SPL sub-bass produces
    measurable skin displacement.  Mechanoreceptors in the skin
    (particularly Pacinian corpuscles, which are most sensitive at
    \SIrange{200}{300}{\hertz} but respond down to
    ${\sim}\SI{20}{\hertz}$) contribute to the tactile sensation of
    bass~\cite{Todd2000}.

  \item \textbf{Bone conduction.}  Sub-bass can excite the skull and
    skeleton directly, stimulating the vestibular system and producing
    a sensation that is perceived as ``visceral'' even though it does
    not involve abdominal resonance.
\end{enumerate}

\subsection{Implications for concert sound design}
\label{sec:concert-design}

Sound engineers have long known empirically that sub-bass ``impact''
requires physical coupling---hence the prevalence of tactile
transducers (``butt-kickers'') in cinema seats and the careful
placement of sub-woofer arrays to maximise floor
vibration~\cite{Leventhall2007}.  Our results provide quantitative
support for this practice: airborne sub-bass alone, even at the
extreme SPLs produced by modern systems, can only produce
$\sim$\SI{0.35}{\percent} of the perception threshold.  Effective
bass reproduction is fundamentally a \emph{structural vibration}
problem.

This has practical implications for venue design.  Venues with
vibration-isolating floors (e.g., floating concrete slabs) may
inadvertently reduce the perceived ``punch'' of sub-bass, even if the
acoustic field is unchanged.  Conversely, lightweight floors that
transmit vibration efficiently will enhance the visceral experience.

\subsection{Pipe-organ acoustics and pew resonance}
\label{sec:pipe-organ}

The pipe organ is a unique case because it produces extreme sub-bass
levels (\SI{90}{\decibel} to \SI{105}{\decibel} at
\SIrange{16}{50}{\hertz}) in enclosed stone buildings where floor
vibration might be expected to be minimal.  However, the wooden pews
typical of churches provide a potential mechanical coupling path.

Using Euler--Bernoulli beam theory for a simply supported oak pew seat
(\SI{2.5}{\metre} span, \SI{40}{\milli\metre} thickness,
$E = \SI{11}{\giga\pascal}$, $\rho = \SI{600}{\kilogram\per\metre\cubed}$),
the first bending resonance is $f_1 \approx \SI{12}{\hertz}$, with
$f_2 \approx \SI{50}{\hertz}$.  Shorter pews (\SI{2.0}{\metre}) have
$f_1 \approx \SI{19}{\hertz}$, placing the first mode directly in the
sub-bass band.

A 32-foot organ pipe's fundamental at \SI{16}{\hertz} with harmonics at
\SI{32}{\hertz} and \SI{48}{\hertz} can therefore excite pew bending
modes, converting the airborne sound field into mechanical vibration of
the seat.  This provides a structure-borne pathway even in stone
buildings with rigid floors, potentially explaining the visceral
sensation reported by congregants during full-organ passages.

Quantitative validation would require in-situ measurements of pew
vibration during organ performance, which we identify as a valuable
direction for future experimental work.

\subsection{Near-field considerations}
\label{sec:near-field-discussion}

The near-field pressure enhancement of $+\SI{6}{}$--\SI{10}{\decibel}$
at distances of \SIrange{0.5}{1.5}{\metre} from a subwoofer stack is
significant for front-row concert audiences.  At \SI{0.5}{\metre}, the
effective SPL can reach \SI{124}{\decibel} at \SI{40}{\hertz},
increasing the airborne displacement to ${\sim}\SI{1}{\percent}$ of the
perception threshold.  Combined with the structure-borne pathway, the
total displacement in the near field may approach or exceed
perceptibility.

This has implications for concert safety: audience members standing
directly in front of subwoofer stacks experience substantially higher
effective pressures than the far-field SPL would suggest.

\subsection{Occupational exposure considerations}
\label{sec:occupational}

For occupational exposure (sound engineers, musicians, nightclub staff),
the relevant regulatory framework is ISO~2631-1, which addresses
whole-body vibration.  Our results confirm that the \emph{acoustic}
component of sub-bass exposure does not contribute meaningfully to
vibration-induced health effects.  However, the structural vibration
component---floor and seat vibration excited by the same
sub-woofers---is a legitimate occupational health concern and should be
assessed separately.

\subsection{Limitations}
\label{sec:limitations}

Several simplifications constrain the present analysis:
%
\begin{enumerate}
  \item The spheroidal shell model assumes a uniform, isotropic wall.
    The real abdominal wall is layered and anisotropic; however, these
    refinements primarily affect the resonance frequency and quality
    factor, not the fundamental coupling physics.

  \item The energy-consistent displacement formulation assumes
    steady-state conditions.  Transient effects (e.g., the EDM ``drop'')
    may briefly exceed steady-state predictions by up to $Q = 4$ times.

  \item We consider only flexural modes.  The breathing mode
    ($n = 0$, $f \approx \SI{2490}{\hertz}$) is driven efficiently by
    uniform pressure but lies far above the sub-bass band.

  \item Concert SPL spectra are taken from published measurements that
    vary significantly between venues and events.  We use representative
    values; actual exposure will differ.

  \item The floor vibration model is a simplified order-of-magnitude
    estimate.  Actual floor response depends on construction type,
    sub-woofer placement, and room modes.

  \item The ISO~2631-1 perception thresholds are conservative for our
    comparison because they describe whole-body perception via the
    mechanical pathway.  True perception thresholds for airborne-induced
    abdominal displacement may be even higher.

  \item We do not model the chest wall, skin, or vestibular system,
    which may contribute significantly to bass perception.
\end{enumerate}
